\section{積分を含む等式から関数を決める}\label{節:積分方程式}
未知の関数が積分の中にも現れる等式を，\textbf{積分方程式}という。上端が変数なら，微分によって積分を外せる場合がある。ただし，微分は定数の情報を失うため，元の等式から得られる値の条件を残し，最後に元の式へ代入して確かめる。

\begin{example}{微分と初期値から関数を求める}
連続関数 $f$ が，すべての実数 $x$ について
\[
 f(x)=2+\int_0^x f(t)\dd t
\]
を満たすとき，$f$ を求めよ。
\end{example}
$x=0$ を代入すると，まず $f(0)=2$ と分かる。$f$ が連続なので，右辺の積分は基本定理により微分可能である。したがって左辺の $f$ も微分可能となり，両辺を微分して $f'(x)=f(x)$ を得る。

この等式だけで直ちに $f=\e^x$ としてはいけない。定数倍も同じ微分の関係を満たすからである。$\e^{-x}$ を掛けた関数を微分すると，
\[
 \frac{\dd}{\dd x}\{\e^{-x}f(x)\}
 =\e^{-x}\{f'(x)-f(x)\}=0
\]
となる。導関数が $0$ なので $\e^{-x}f(x)$ は定数であり，$x=0$ の値からその定数は $2$ である。よって候補は $f(x)=2\e^x$ となる。

最後に元の式を確認する。
\[
 2+\int_0^x2\e^t\dd t
 =2+[2\e^t]_0^x=2+2\e^x-2=2\e^x.
\]
確かに成り立つので，求める関数は $f(x)=2\e^x$ である。導いた条件が必要条件であるだけでなく，元の積分方程式を満たすことまで確かめた。

ここまでで，定積分の意味と，その計算を微分に結び付ける根拠がそろった。次章からは，原始関数を見つけるための式の整理を詳しく扱う。どの公式を使う場合も，求めた原始関数を微分して被積分関数に戻るかを確認できる。複雑な計算ほど，この戻り道が役に立つ。

